% !PRINTED 621
\weberpage{621}
% Faint lines above §181 are reversed verso show-through; no direct recto body text is present.

\begin{center}
\S.\,181.

\gesperrt{Anwendung auf die metacyklischen Gleichungen}
$5^{\text{ten}}$ \gesperrt{Grades.}
\end{center}

\quad Wir machen noch eine Anwendung der allgemeinen Sätze
auf die Gleichungen $5^{\text{ten}}$ Grades. Ist $n = 5$, so hat die cyklische
Gruppe $5$, die halbmetacyklische $10$ und die volle metacyklische
Gruppe $20$ Permutationen. Eine metacyklische Function genügt
bei einer allgemeinen Gleichung $5^{\text{ten}}$ Grades einer Resolvente
$6^{\text{ten}}$ Grades, und wenn diese Gleichung $6^{\text{ten}}$ Grades eine rationale
Wurzel hat, so ist die gegebene Gleichung $5^{\text{ten}}$ Grades metacyklisch.
Die halbmetacyklischen Functionen, d.\,h. die Functionen,
die die Permutationen der halbmetacyklischen Gruppe
gestatten, genügen einer Resolvente $12^{\text{ten}}$ Grades, die aber nach
Adjunction der Quadratwurzel aus der Discriminante in zwei
Factoren $6^{\text{ten}}$ Grades zerfällt.

\quad Wenn man die erzeugenden Substitutionen für den Modul $5$
\[
s = (z, z+1),\quad t = (z, 2z),\quad t^2 = (z, 4z)
\]
auf die Ziffern $0, 1, 2, 3, 4$ anwendet, so ergeben sich die Vertauschungen
\weq{1}{
\begin{array}{lccccc}
 & 0, & 1, & 2, & 3, & 4\\
(s) & 1, & 2, & 3, & 4, & 0\\
(t) & 0, & 2, & 4, & 1, & 3\\
(t^2) & 0, & 4, & 3, & 2, & 1.
\end{array}
}

\quad Wenden wir die Substitutionen $s, t^2, t$ auf die Paare von
Ziffern $(0,1), (1,2), (2,3), (3,4), (4,0)$ an, so folgt
\weq{2}{
\begin{array}{lccccc}
 & (0,1), & (1,2), & (2,3), & (3,4), & (4,0)\\
(s) & (1,2), & (2,3), & (3,4), & (4,0), & (0,1)\\
(t^2) & (0,4), & (4,3), & (3,2), & (2,1), & (1,0);
\end{array}
}
diese Paare bleiben also in ihrer Gesammtheit durch $s$ und
durch $t^2$, und folglich durch die halbmetacyklische Gruppe ungeändert,
und eine symmetrische Function der entsprechenden
Wurzelpaare ist halbmetacyklisch. Die Substitution $t$ bewirkt
die Vertauschungen
\weq{3}{
\begin{array}{lccccc}
 & (0,1), & (1,2), & (2,3), & (3,4), & (4,0)\\
(t) & (0,2), & (2,4), & (4,1), & (1,3), & (3,0),
\end{array}
}
%CONT>>
